\renewcommand{\Anno}{2013}

\renewcommand{\NumeroProblema}{3}

\renewcommand{\Testo}{
Consideriamo un rombo inscritto in un rettangolo, cioè un rombo i cui 4 vertici stanno sui lati del rettangolo.
\begin{enumerate}
	\item[(a)] Supponiamo che ogni lato del rettangolo contenga (all'interno o agli estremi) almeno un vertice del rombo. Dimostrare che in tale ipotesi il centro del rombo coincide con il centro del rettangolo (il centro di un rombo o di un rettangolo è l'intersezione delle diagonali).

	\item[(b)] Determinare in quali rettangoli esistono rombi inscritti il cui centro non coincide con il centro del rettangolo.

	\item[(c)] Indicate con $a$ e $b$ le lunghezze dei lati del rettangolo, con $a\geq b$, determinare i possibili valori della lunghezza del lato di un qualunque rombo inscritto.
\end{enumerate}
}

\renewcommand{\Risposta}{
\begin{enumerate}
	\item[(a)] \ldots

	\item[(b)] In tutti i rettangoli che non sono quadrati.

	\item[(c)] Indicata con $\ell$ la lunghezza del lato del rombo, i valori possibili di $\ell$ sono tutti e soli quelli tali che
	$$\min\left\{b,\frac{\sqrt{a^{2}+b^{2}}}{2}\right\}
	\leq\ell\leq\frac{a^{2}+b^{2}}{2a}.$$

\end{enumerate}
}

\renewcommand{\Aiutino}{
\begin{enumerate}

	\item[(a)] Dimostrare che i 4 triangolini esterni al rombo sono a 2 a 2 congruenti.

	\item[(b)] Basta pensare ai quadrati inscritti.

	\item[(c)] Dai punti precedenti sono emerse due categorie di rombi inscritti.  Per quelli con il centro coincidente con il centro del rettangolo conviene esprimere la lunghezza del lato in funzione dell'angolo che le diagonali del rombo formano con i lati del rettangolo.  Per gli altri si tratta di trovare una stima dall'alto ed una dal basso per la lunghezza del lato.
\end{enumerate}
}

\renewcommand{\SoluzioneUfficiale}{\TitoloSoluzione

\DomandaA

Con riferimento alla Figura~1a, mostriamo che i triangoli $APS$ e $CRQ$ sono congruenti.  Infatti $PS=RQ$ in quanto lati dello stesso rombo, mentre $\angle APS=\angle CRQ$ e $\angle ASP=\angle CQR$ per via del parallelismo tra le rette $AB$ e $CD$ ed il parallelismo tra le rette $SP$ ed $RQ$. Questo mostra in particolare che $AP=CR$.

Allo stesso modo si potrebbe dimostrare che i triangoli $BQP$ e $DSR$ sono congruenti.

Ora, con riferimento alla Figura~1b, sia $O$ il punto di intersezione tra le rette $AC$ e $PR$.  Mostriamo che i triangoli $OAP$ e $OCR$ sono congruenti.  Infatti essi hanno i lati $AP$ e $CR$ congruenti, come appena dimostrato, ed inoltre $\angle OAP=\angle OCR$ e $\angle OPA=\angle ORC$ per via del parallelismo tra le rette $AB$ e $CD$.  Ne segue in particolare che $OA=OC$ e $OP=OR$, dunque il punto $O$ è il punto medio sia della diagonale $AC$ del rettangolo, sia della diagonale $PR$ del rombo. Ne segue che il punto $O$ è contemporaneamente il centro del rettangolo e del rombo.

\hfill
\psset{unit=9ex}
\pspicture(-0.5,-1)(3.5,2.5)
\pstGeonode[PosAngle={-135,-45,45,135}](0,0){A}(3,0){B}(3,2){C}(0,2){D}
\psline(A)(B)(C)(D)(A)
\pstInterLL[PointName=none,PointSymbol=none]{A}{C}{B}{D}{O}
\pstMiddleAB[PointName=none,PointSymbol=none]{A}{D}{S'}
\pstRotation[RotAngle=20,PointName=none,PointSymbol=none]{O}{S'}[S'']
\pstInterLL[PosAngle=180]{A}{D}{O}{S''}{S}
\pstInterLL{B}{C}{O}{S''}{Q}
\pstRotation[RotAngle=90,PointName=none,PointSymbol=none]{O}{S''}[P'']
\pstInterLL[PosAngle=270]{A}{B}{O}{P''}{P}
\pstInterLL[PosAngle=90]{D}{C}{O}{P''}{R}
\psline(P)(Q)(R)(S)(P)
\pstMarkAngle[linecolor=magenta,MarkAngleRadius=0.4]{S}{P}{A}{}
\pstMarkAngle[linecolor=magenta,MarkAngleRadius=0.5]{S}{P}{A}{}
\pstMarkAngle[linecolor=magenta,MarkAngleRadius=0.2]{A}{S}{P}{}
\pstMarkAngle[linecolor=magenta,MarkAngleRadius=0.4]{Q}{R}{C}{}
\pstMarkAngle[linecolor=magenta,MarkAngleRadius=0.5]{Q}{R}{C}{}
\pstMarkAngle[linecolor=magenta,MarkAngleRadius=0.2]{C}{Q}{R}{}
\rput(1.5,-0.7){Figura 1a}
\endpspicture
\hfill
\pspicture(-0.5,-1)(3.5,2.5)
\pstGeonode[PosAngle={-135,-45,45,135}](0,0){A}(3,0){B}(3,2){C}(0,2){D}
\psline(A)(B)(C)(D)(A)
\pstInterLL{A}{C}{B}{D}{O}
\pstMiddleAB[PointName=none,PointSymbol=none]{A}{D}{S'}
\pstRotation[RotAngle=20,PointName=none,PointSymbol=none]{O}{S'}[S'']
\pstInterLL[PosAngle=180]{A}{D}{O}{S''}{S}
\pstInterLL{B}{C}{O}{S''}{Q}
\pstRotation[RotAngle=90,PointName=none,PointSymbol=none]{O}{S''}[P'']
\pstInterLL[PosAngle=270]{A}{B}{O}{P''}{P}
\pstInterLL[PosAngle=90]{D}{C}{O}{P''}{R}
\psline[linecolor=cyan](P)(Q)(R)(S)(P)
\psline(R)(P)
\psline(A)(C)
\pstMarkAngle[linecolor=magenta,MarkAngleRadius=0.4]{P}{A}{O}{}
\pstMarkAngle[linecolor=magenta,MarkAngleRadius=0.35]{P}{A}{O}{}
\pstMarkAngle[linecolor=magenta,MarkAngleRadius=0.2]{O}{P}{A}{}
\pstMarkAngle[linecolor=magenta,MarkAngleRadius=0.4]{R}{C}{O}{}
\pstMarkAngle[linecolor=magenta,MarkAngleRadius=0.35]{R}{C}{O}{}
\pstMarkAngle[linecolor=magenta,MarkAngleRadius=0.2]{O}{R}{C}{}
\rput(1.5,-0.7){Figura 1b}
\endpspicture
\hfill\mbox{}

\DomandaB

Se il rettangolo non è un quadrato, cioè ha i lati di lunghezza diversa, allora esistono sempre dei rombi inscritti il cui centro non coincide con il centro del rettangolo.  Basta infatti considerare i quadrati inscritti con i lati paralleli ai lati del rettangolo.  Uno di tali quadrati ha il centro coincidente con il centro del rettangolo (vedi Figura~2a), ma basta traslare il quadrato per fare in modo che il centro non coincida più con quello del rettangolo (vedi Figura~2b e Figura~2c).\bigskip

{\centering
\hfill
\psset{unit=6ex}
\pspicture(0,-1)(3,2)
\pstGeonode[PointName=none](0,0){A}(3,0){B}(3,2){C}(0,2){D}
\pstGeonode[PointName=none](0.5,0){P}(2.5,0){Q}(2.5,2){R}(0.5,2){S}
\psline(A)(B)(C)(D)(A)
\psline[linecolor=cyan](P)(Q)(R)(S)(P)
\psline[linecolor=magenta,linestyle=dashed](A)(C)
\psline[linecolor=magenta,linestyle=dashed](B)(D)
\psline[linecolor=blue,linestyle=dashed](P)(R)
\psline[linecolor=blue,linestyle=dashed](Q)(S)
\pstInterLL[PointName=none]{A}{C}{B}{D}{O}
\rput(1.5,-0.7){Figura 2a}
\endpspicture
\hfill
\pspicture(0,-1)(3,2)
\pstGeonode[PointName=none](0,0){A}(3,0){B}(3,2){C}(0,2){D}
\pstGeonode[PointName=none](0.8,0){P}(2.8,0){Q}(2.8,2){R}(0.8,2){S}
\psline(A)(B)(C)(D)(A)
\psline[linecolor=cyan](P)(Q)(R)(S)(P)
\psline[linecolor=magenta,linestyle=dashed](A)(C)
\psline[linecolor=magenta,linestyle=dashed](B)(D)
\psline[linecolor=blue,linestyle=dashed](P)(R)
\psline[linecolor=blue,linestyle=dashed](Q)(S)
\pstInterLL[PointName=none]{A}{C}{B}{D}{O}
\pstInterLL[PointName=none]{P}{R}{Q}{S}{O'}
\rput(1.5,-0.7){Figura 2b}
\endpspicture
\hfill
\pspicture(0,-1)(3,2)
\pstGeonode[PointName=none](0,0){A}(3,0){B}(3,2){C}(0,2){D}
\pstGeonode[PointName=none](0.3,0){P}(2.3,0){Q}(2.3,2){R}(0.3,2){S}
\psline(A)(B)(C)(D)(A)
\psline[linecolor=cyan](P)(Q)(R)(S)(P)
\psline[linecolor=magenta,linestyle=dashed](A)(C)
\psline[linecolor=magenta,linestyle=dashed](B)(D)
\psline[linecolor=blue,linestyle=dashed](P)(R)
\psline[linecolor=blue,linestyle=dashed](Q)(S)
\pstInterLL[PointName=none]{A}{C}{B}{D}{O}
\pstInterLL[PointName=none]{P}{R}{Q}{S}{O'}
\rput(1.5,-0.7){Figura 2c}
\endpspicture
\hfill
\mbox{}}

Se invece il rettangolo è un quadrato, allora tutti i rombi inscritti hanno il centro coincidente con il centro del quadrato.  Se infatti così non fosse, dal punto~(a) sapremmo che esiste un lato del quadrato che non contiene vertici del rombo, e di conseguenza esisterebbe un lato del quadrato che contiene almeno due vertici del rombo, non entrambi agli estremi.

Supponiamo dunque, senza perdita di generalità, che il lato $AB$ del quadrato contenga i vertici $P$ e $Q$ del rombo (si tratta di vertici consecutivi, ovviamente, come indicato nella Figura~3a).

Consideriamo allora i restanti vertici $R$ ed $S$ del rombo. Poiché il lato $RS$ del rombo deve essere parallelo al lato $PQ$, ci sono solo due possibilità.
\begin{itemize}
	\item Possibilità 1: $R$ sta sul lato $BC$ ed $S$ sta sul lato $DA$ (vedi Figura~3b).  Questo non è possibile perché si avrebbe che
	$$AB>PQ=RS= AB.$$

	\item Possibilità 2: $R$ ed $S$ stanno entrambi sul lato $CD$ del quadrato (vedi Figura~3c).  Anche questo però non è possibile in quanto si avrebbe che
	$$PQ=QR\geq BC=AB>PQ.$$
\end{itemize}

{\centering
\hfill
\psset{unit=7ex}
\pspicture(0,-1)(2,2.5)
\pstGeonode[PosAngle={-135,-45,45,135}](0,0){A}(2,0){B}(2,2){C}(0,2){D}
\pstGeonode[PosAngle={270,270}](0.4,0){P}(1.8,0){Q}
\psline(A)(B)(C)(D)(A)
\rput(1,-0.9){Figura 3a}
\endpspicture
\hfill
\pspicture(0,-1)(2,2.5)
\pstGeonode[PosAngle={-135,-45,45,135}](0,0){A}(2,0){B}(2,2){C}(0,2){D}
\pstGeonode[PosAngle={270,270,0,180}](0.4,0){P}(1.8,0){Q}(2,1.5){R}(0,1.5){S}
\psline(A)(B)(C)(D)(A)
\psline[linecolor=cyan](P)(Q)(R)(S)(P)
\rput(1,-0.9){Figura 3b}
\endpspicture
\hfill
\pspicture(0,-1)(2,2.5)
\pstGeonode[PosAngle={-135,-45,45,135}](0,0){A}(2,0){B}(2,2){C}(0,2){D}
\pstGeonode[PosAngle={270,270,90}](0.4,0){P}(1.8,0){Q}(1.7,2){R}(0.3,2){S}
\psline(A)(B)(C)(D)(A)
\psline[linecolor=cyan](P)(Q)(R)(S)(P)
\rput(1,-0.9){Figura 3c}
\endpspicture
\hfill
\mbox{}}

\DomandaC

Dalle discussioni svolte ai punti precedenti, sappiamo che in un generico rettangolo ci sono due tipi di rombi inscritti: quelli per cui ogni lato del rettangolo contiene un vertice del rombo e quelli con due vertici su uno stesso lato del rettangolo.  Dai punti precedenti sappiamo che i primi hanno il centro coincidente con il centro del rettangolo, mentre i secondi hanno necessariamente due vertici su un lato lungo del rettangolo e due vertici sul lato opposto.

Mostreremo ora che per i primi la lunghezza $\ell$ del lato varia tra il minimo assunto quando i vertici del rombo sono i punti medi dei lati del rettangolo (vedi Figura~4a) ed il massimo assunto quando due vertici del rombo coincidono con i vertici del rettangolo (vedi Figura~4b).  Per i rombi del secondo tipo, la lunghezza $\ell$ del lato varia tra il minimo assunto quando il rombo è uno dei quadrati inscritti con lati paralleli ai lati del rettangolo (vedi Figura~4c) ed il massimo assunto ancora una volta nella configurazione della Figura~4b.  Il rombo della configurazione della Figura~4b appartiene ad entrambe le tipologie, di cui rappresenta il caso limite.\bigskip

{\centering
\hfill
\psset{unit=7ex}
\pspicture(0,-1)(3,2)
\pstGeonode[PointName=none](0,0){A}(3,0){B}(3,2){C}(0,2){D}
\pstMiddleAB[PointName=none]{A}{B}{P}
\pstMiddleAB[PointName=none]{B}{C}{Q}
\pstMiddleAB[PointName=none]{C}{D}{R}
\pstMiddleAB[PointName=none]{D}{A}{S}
\psline(A)(B)(C)(D)(A)
\psline[linecolor=cyan](P)(Q)(R)(S)(P)
\rput(1.5,-0.5){Figura 4a}
\endpspicture
\hfill
\pspicture(0,-1)(3,2)
\pstGeonode[PointName=none](0,0){A}(3,0){B}(3,2){C}(0,2){D}
\pstInterLL[PointName=none,PointSymbol=none]{A}{C}{B}{D}{O}
\pstRotation[RotAngle=90,PointName=none,PointSymbol=none]{O}{A}[A']
\pstInterLL[PointName=none]{A}{B}{O}{A'}{Q}
\pstInterLL[PointName=none]{C}{D}{O}{A'}{S}
\psline(A)(B)(C)(D)(A)
\psline[linecolor=cyan](A)(Q)(C)(S)(A)
\rput(1.5,-0.5){Figura 4b}
\endpspicture
\hfill
\pspicture(0,-1)(3,2)
\pstGeonode[PointName=none](0,0){A}(3,0){B}(3,2){C}(0,2){D}
\pstGeonode[PointName=none](0.3,0){P}(2.3,0){Q}(2.3,2){R}(0.3,2){S}
\psline(A)(B)(C)(D)(A)
\psline[linecolor=cyan](P)(Q)(R)(S)(P)
\rput(1.5,-0.5){Figura 4c}
\endpspicture
\hfill
\mbox{}}

Ora è facile vedere che la lunghezza del lato del rombo della Figura~4a è metà della diagonale del rettangolo, cioè $\sqrt{a^{2}+b^{2}}/2$, mentre è banalmente $b$ nel caso della Figura~4c.  Mostreremo che nel caso della Figura~4b la lunghezza del lato è $(a^{2}+b^{2})/(2a)$, da cui seguirà immediatamente che la lunghezza $\ell$ del lato di ogni rombo inscritto verifica
$$\min\left\{b,\frac{\sqrt{a^{2}+b^{2}}}{2}\right\}
\leq\ell\leq\frac{a^{2}+b^{2}}{2a}.$$

Consideriamo dunque un generico rombo del primo tipo, cioè con un vertice su ogni lato del rettangolo.  Con riferimento alla Figura~5a, indichiamo con $\theta$ l'angolo che le diagonali del rombo formano con i lati del rettangolo.  Detto $O$ il centro del rettangolo (e del rombo), e detti $M$ e $K$ i punti medi dei lati $AB$ ed $AD$, rispettivamente, con un po' di trigonometria elementare si ha che
$$OS=\frac{OK}{\cos\theta}=\frac{a}{2\cos\theta},
\hspace{3em}
OP=\frac{OM}{\cos\theta}=\frac{b}{2\cos\theta},$$
da cui
\begin{equation}
	\ell=\sqrt{OS^{2}+OP^{2}}=\frac{\sqrt{a^{2}+b^{2}}}{2\cos\theta}.
	\label{rombi:ell}
\end{equation}

{\centering
\hfill
\psset{unit=9ex}
\pspicture(-0.3,-1)(3.3,2.5)
\pstGeonode[PosAngle={-135,-45,45,135}](0,0){A}(3,0){B}(3,2){C}(0,2){D}
\pstMiddleAB[PosAngle=270]{A}{B}{M}
\pstMiddleAB{A}{D}{K}
\psline(A)(B)(C)(D)(A)
\pstInterLL[PosAngle=-20]{A}{C}{B}{D}{O}
\pstMiddleAB[PointName=none,PointSymbol=none]{A}{D}{S'}
\pstRotation[RotAngle=20,PointName=none,PointSymbol=none]{O}{S'}[S'']
\pstInterLL[PosAngle=180]{A}{D}{O}{S''}{S}
\pstInterLL{B}{C}{O}{S''}{Q}
\pstRotation[RotAngle=90,PointName=none,PointSymbol=none]{O}{S''}[P'']
\pstInterLL[PosAngle=270]{A}{B}{O}{P''}{P}
\pstInterLL[PosAngle=90]{D}{C}{O}{P''}{R}
\psline[linecolor=cyan](P)(Q)(R)(S)(P)
\psline[linecolor=blue,linestyle=dashed](S)(Q)
\psline[linecolor=blue,linestyle=dashed](P)(R)
\psline[linecolor=blue,linestyle=dashed](K)(O)(M)
\pstMarkAngle[linecolor=magenta,MarkAngleRadius=0.35,LabelSep=0.7]{K}{O}{S}{$\theta$}
\pstMarkAngle[linecolor=magenta,MarkAngleRadius=0.3]{K}{O}{S}{}
\pstMarkAngle[linecolor=magenta,MarkAngleRadius=0.35,LabelSep=0.6]{M}{O}{P}{$\theta$}
\pstMarkAngle[linecolor=magenta,MarkAngleRadius=0.3]{M}{O}{P}{}
\rput(1.5,-0.7){Figura 5a}
\endpspicture
\hfill
\pspicture(-0.3,-1)(3.8,2.5)
\pstGeonode[PosAngle={-135,-45,45,135}](0,0){A}(3.5,0){B}(3.5,2){C}(0,2){D}
\psline(A)(B)(C)(D)(A)
\pstGeonode[PosAngle=270](0.2,0){P}(2.3,0){Q}
\pstInterLC[PosAngle=90,PointNameA=none,PointSymbolA=none]{D}{C}{P}{Q}{S'}{S}
\pstTranslation[PosAngle=90]{P}{Q}{S}[R]
\psline[linecolor=cyan](P)(Q)(R)(S)(P)
\pstProjection{A}{B}{R}[R']
\psline[linecolor=magenta,linestyle=dashed](R)(R')
\rput(1.75,-0.7){Figura 5b}
\endpspicture
\hfill
\mbox{}}

È pertanto evidente che $\ell$ assume il valore minimo quando $\cos\theta$ è massimo, cioè quando $\theta=0$, cioè nella configurazione della Figura~4a. Analogamente, $\ell$ assume il valore massimo quando $\cos\theta$ è il minimo possibile, cioè quando $\theta$ è il massimo possibile, il che accade quando $S$ coincide con $A$. Siamo dunque nella configurazione della Figura~4b, in cui
$$\cos\theta=\frac{OK}{OA}=\frac{a}{\sqrt{a^{2}+b^{2}}}.$$

Sostituendo tale valore nella (\ref{rombi:ell}) si ottiene
$$\ell=\frac{a^{2}+b^{2}}{2a},$$
che è la lunghezza del lato del rombo nella configurazione della Figura~4b (che già sappiamo essere la massima lunghezza possibile).

Potendo $\theta$ assumere tutti i valori intermedi, è evidente che anche $\ell$ potrà assumere tutti i valori intermedi.

Consideriamo infine un generico rombo del secondo tipo, cioè con i vertici sui lati lunghi del rettangolo. Con riferimento alla Figura~5b, indichiamo con $R'$ la proiezione di $R$ su $AB$. È ora evidente che
$$\ell=RQ\geq RR'=CB=b,$$
da cui $\ell\geq b$, con uguaglianza se e solo se $RQ$ è parallelo a $CB$, cioè se e solo se il rombo è un quadrato con i lati paralleli a quelli del rettangolo. Inoltre
$$\ell+\sqrt{\ell^{2}-b^{2}}=PQ+QR'\leq AB=a.$$

Si ottiene così la disequazione $\ell+\sqrt{\ell^{2}-b^{2}}\leq a$, la quale (limitandosi ai valori $\ell\geq b$) ha come soluzione
$$\ell\leq\frac{a^{2}+b^{2}}{2a}.$$

L'uguaglianza si ha proprio quando $PQ+QR'= AB$, cioè nella configurazione della Figura~4b. Non è difficile mostrare che tutti i valori intermedi di $\ell$ sono effettivamente realizzabili: basta prendere $P$ coincidente con $A$, tracciare la circonferenza con centro in $P$ e raggio $\ell$, definire $Q$ ed $S$ come le intersezioni tra la circonferenza ed i lati $AB$ e $CD$, ed infine definire $R$ ``completando il parallelogrammo''. Le restrizioni su $\ell$ assicurano proprio che $Q$ ed $S$ esistono, e che $R$ sta ancora sul lato $CD$ del rettangolo.
}

\renewcommand{\SoluzioniAlternative}{
Presentiamo una soluzione alternativa per il punto~(a). Con riferimento alla Figura~1a, consideriamo un sistema cartesiano con l'origine in $A$ e gli assi paralleli ai lati del rettangolo. Dette $a$ e $b$ le lunghezze dei lati del rettangolo, le coordinate dei vertici del rettangolo saranno quindi
$$A=(0,0),
\hspace{3em}
B=(a,0),
\hspace{3em}
C=(a,b),
\hspace{3em}
D=(0,b).$$

Le coordinate dei vertici del rombo saranno invece
$$P=(p,0),
\hspace{3em}
Q=(a,q),
\hspace{3em}
R=(r,b),
\hspace{3em}
S=(0,s)$$
per opportuni valori dei parametri $p$, $q$, $r$, $s$. Calcoliamo ora le coordinate del centro $O$ del rombo. Essendo $O$ il punto medio del segmento $SQ$, la sua ascissa sarà la media delle ascisse di $S$ e $Q$, quindi $a/2$. Essendo $O$ il punto medio del segmento $PR$, la sua ordinata sarà la media delle ordinate di $P$ ed $R$, quindi $b/2$. Ne segue che $O=(a/2,b/2)$, che coincide quindi con il centro del rettangolo.

Si noti che il trucco consiste nell'usare il segmento $SQ$ per il calcolo dell'ascissa ed il segmento $PR$ per il calcolo dell'ordinata del centro.  In questo modo i parametri $p$, $q$, $r$, $s$ non entrano in gioco in questa prima fase.  A posteriori, usando ora il segmento $SQ$ per il calcolo dell'ordinata ed il segmento $PR$ per il calcolo dell'ascissa, si ottiene che $s+q=b$ e $p+r=a$, che è proprio la simmetria alla base della soluzione per via sintetica.



\bigskip

Segnaliamo ora un modo alternativo per trattare nel punto~(c) il caso dei rombi con un vertice su ogni lato del rettangolo. Con riferimento alla Figura~6a, poniamo $AP=x$ e $AS=y$. Grazie alla simmetria dimostrata al punto~(a), tutte le restanti lunghezze risulteranno determinate come in figura. In particolare avremo che
$$PS^{2}=\ell^{2}=x^{2}+y^{2},$$
$$PQ^{2}=\ell^{2}=(a-x)^{2}+(b-y)^{2},$$
il che porta al sistema
$$\left\{
\begin{array}{l}
	x^{2}+y^{2}=\ell^{2},  \\
	x^{2}+y^{2}-2ax-2by+a^{2}+b^{2}=\ell^{2}.
\end{array}
\right.$$

Sostituendo la prima equazione nella seconda, il sistema diventa
$$\left\{
\begin{array}{l}
	x^{2}+y^{2}=\ell^{2},  \\
	2ax+2by=a^{2}+b^{2}.
\end{array}
\right.$$

{\centering
\hfill
\psset{unit=9ex}
\pspicture(-0.5,-1)(3.5,2.5)
\pstGeonode[PosAngle={-135,-45,45,135}](0,0){A}(3,0){B}(3,2){C}(0,2){D}
\psline(A)(B)(C)(D)(A)
\pstInterLL[PointName=none,PointSymbol=none]{A}{C}{B}{D}{O}
\pstMiddleAB[PointName=none,PointSymbol=none]{A}{D}{S'}
\pstRotation[RotAngle=15,PointName=none,PointSymbol=none]{O}{S'}[S'']
\pstInterLL[PosAngle=180]{A}{D}{O}{S''}{S}
\pstInterLL{B}{C}{O}{S''}{Q}
\pstRotation[RotAngle=90,PointName=none,PointSymbol=none]{O}{S''}[P'']
\pstInterLL[PosAngle=270]{A}{B}{O}{P''}{P}
\pstInterLL[PosAngle=90]{D}{C}{O}{P''}{R}
\psline[linecolor=cyan](P)(Q)(R)(S)(P)
\pcline[linestyle=none](A)(P)
\bput{:U}{\magenta $x$}
\pcline[linestyle=none](P)(B)
\bput{:U}{\magenta $a-x$}
\pcline[linestyle=none](D)(R)
\aput{:U}{\magenta $a-x$}
\pcline[linestyle=none](R)(C)
\aput{:U}{\magenta $x$}
\pcline[linestyle=none](A)(S)
\aput{0}{\magenta $y$}
\pcline[linestyle=none](S)(D)
\aput{0}{\magenta $b-y$}
\pcline[linestyle=none](B)(Q)
\bput{0}{\magenta $b-y$}
\pcline[linestyle=none](Q)(C)
\bput{0}{\magenta $y$}
\pcline[linestyle=none](S)(P)
\aput{0}{\magenta $\ell$}
\pcline[linestyle=none](P)(Q)
\aput{0}{\magenta $\ell$}
\rput(1.5,-0.7){Figura 6a}
\endpspicture
\hfill\hfill
\pspicture(-0.5,-1)(2.5,2.5)
\psline[linewidth=0.7\pslinewidth,linecolor=cyan]{->}(-0.5,0)(2.5,0)
\psline[linewidth=0.7\pslinewidth,linecolor=cyan]{->}(0,-0.5)(0,2.3)
\pstGeonode[PosAngle={225,180,230}](0,0){O}(0,2){B'}(1.5,0){A'}
\pstProjection[PosAngle=200]{A'}{B'}{O}[K]
\pstSymO[PosAngle=200]{K}{A'}[C]
\psclip{\psframe[linestyle=none](-0.5,-0.5)(2.5,2.3)}
\pstLineAB[nodesep=-1]{A'}{B'}
\pstCircleOA[linecolor=magenta]{O}{K}
\pstCircleOA[linecolor=magenta]{O}{A'}
\pscircle[linecolor=magenta,linestyle=dashed](O){1.7}
\endpsclip

\rput(1,-0.7){Figura 6b}
\endpspicture
\hfill
\mbox{}}

Ora si tratta di capire per quali valori del parametro $\ell$ il sistema ammette una soluzione $(x,y)$ con $0\leq x\leq a$ e $0\leq y\leq b$.  Nel piano cartesiano (vedi Figura~6b), la prima equazione rappresenta la circonferenza con centro nell'origine e raggio $\ell$, la seconda equazione rappresenta la retta passante per i punti
$$A'=\left(\frac{a^{2}+b^{2}}{2a},0\right), \hspace{4em}
B'=\left(0,\frac{a^{2}+b^{2}}{2b}\right).$$

Si noti che, tenendo conto dell'ipotesi $a\geq b$, la distanza di $B'$ dall'origine è maggiore od uguale della distanza di $A'$ dall'origine.

Affinché ci sia intersezione, $\ell$ deve essere maggiore o uguale alla distanza dell'origine dalla retta. Tale distanza si calcola con una nota formula e ci permette di concludere che
$$\ell\geq\frac{\sqrt{a^{2}+b^{2}}}{2}.$$

Quando c'è uguaglianza, la soluzione del sistema è la proiezione dell'origine sulla retta, che si verifica essere il punto $K=(a/2,b/2)$. Questa scelta corrisponde alla configurazione della Figura~4a.

Viceversa, il massimo valore di $\ell$ per cui c'è intersezione con $0\leq x\leq a$ e $0\leq y\leq b$ è quello per cui la circonferenza passa per $A'$. In tal caso si ha
$$\ell=\frac{a^{2}+b^{2}}{2a},$$
ed il sistema ammette 2 soluzioni: $(x,y)=A'$, che corrisponde alla configurazione della Figura~4b, e $(x,y)=C$, con
$$C=\left(\frac{a^{2}-b^{2}}{2a},b\right),$$
che corrisponde al simmetrico della Figura~4b rispetto all'asse verticale del rettangolo. È evidente che per valori maggiori di $\ell$ le soluzioni del sistema hanno $y<0$ oppure $y>b$.
}

\renewcommand{\Osservazioni}{\TitoloOsservazione

Abbiamo dimostrato che il massimo valore possibile per la lunghezza $\ell$ del lato del rombo inscritto si ottiene sempre nella configurazione della Figura~4b. Al contrario, il valore minimo di $\ell$ si ottiene nella configurazione della Figura~4a o della Figura~4c a seconda del rapporto tra $a$ e $b$. In particolare, il minimo è rappresentato dalla  Figura~4c se e solo se $a\geq b\sqrt{3}$, cioè se e solo se il rettangolo è ``sufficientemente allungato''.

}

\renewcommand{\Commenti}{}
